\section{A finite-stencil inverse and its sampling law}
\label{sec:finite-stencil}
The global stopping formula can be evaluated on a fixed finite set of
nominal centers. In particular, neither its spatial aperture nor its
forcing precision needs to grow with an iteration horizon. The finite
inverse is an optimal stopping linear program on a known killed walk.
Its killing boundary is computational; no assertion that the physical
occupation vanishes on that boundary is made.

Throughout this section the configuration and launch law satisfy the
hypotheses of Section~\ref{sec:global-response}. Put
\begin{equation}\label{eq:stencil-constants}
 K=\max\{1,\lceil(D+\Delta)/t\rceil\},\quad
 S=\{-K,\ldots,K\}^2,\quad M=(2K+1)^2,\quad
 H_S=K^2+(K+1)^2.
\end{equation}
For a target $x$, the nominal stencil is $x+tS$. Let $Q$ be the
substochastic compass matrix on $S$: a step leaving $S$ is killed and
its mass is not redistributed. Write $o=(0,0)$, let $e_o$ be its unit
coordinate vector, and set
\begin{equation}\label{eq:stencil-green}
 G=(I-Q)^{-1},\qquad h=G\one,\qquad
 \mathcal M_o=\{m\ge0:(I-Q^{\mathsf T})m\le e_o\}.
\end{equation}
Existence of $G$ and compactness of $\mathcal M_o$ are included in the
next lemma. For a real vector $f$ on $S$ define
\begin{equation}\label{eq:stencil-functional}
 \mathcal V(f)=\max_{m\in\mathcal M_o}(-f^{\mathsf T}m).
\end{equation}
The matrix, feasible set, and all their coefficients are independent
of the obstacles, the launch density, and the data.

\begin{lemma}[The killed stopping polytope]
\label{lem:stencil-polytope}
The killed compass walk exits $S$ almost surely, and
$0<h_z\le H_S$ for every $z\in S$. The set $\mathcal M_o$ is a
nonempty compact polytope. Its elements are precisely the expected
continuation-visit vectors of randomized stopping rules that stop no
later than exit from $S$. Moreover,
\begin{equation}\label{eq:stencil-budget}
 0\le m_z\le G_{oz},\qquad \sum_zm_z\le h_o\le H_S.
\end{equation}
For every real $f$ one has the attained dual formula
\begin{equation}\label{eq:stencil-dual}
 \mathcal V(f)=
 \min\{u_o:u\ge0,\ (I-Q)u\ge-f\}.
\end{equation}
Equivalently, $\mathcal V(f)$ is the value at $o$ of the unique
solution $u=\max\{0,Qu-f\}$.
\end{lemma}
\begin{proof}
Use lattice units and let $\tau_S$ be the first exit time. The process
$|X_n|^2-n$ is a martingale. At $n\wedge\tau_S$ the squared norm is
at most $K^2+(K+1)^2$. Bounded stopping, followed by monotone
convergence, gives
$\E_z\tau_S\le H_S-|z|^2\le H_S$. Thus exit is almost sure,
$Q^n\to0$, and the nonnegative Green series $\sum_{n\ge0}Q^n$
converges to $G$. Its row sums are the expected exit times.

If $m_z$ counts expected continuations at $z$ and $s_z$ counts
expected stops there, conservation of visits gives
\[
                 m+s=e_o+Q^{\mathsf T}m,\qquad m,s\ge0.
\]
Conversely, given these equations, continue at $z$ with probability
$m_z/(m_z+s_z)$, interpreted as zero when the denominator vanishes.
The resulting continuation transition matrix is bounded entrywise by
$Q$ and is therefore transient. Its visit equations have a unique
solution, namely $m+s$, so its continuation and stopping masses are
exactly $m$ and $s$. This proves the asserted realization of every
feasible vector; history-dependent rules satisfy the same conservation
identities.

Multiplication by the nonnegative matrix $G^{\mathsf T}$ yields
\[
             m=G^{\mathsf T}(e_o-s)\le G^{\mathsf T}e_o.
\]
Summation proves \eqref{eq:stencil-budget}, which also proves
compactness. The zero vector is feasible. The dual feasible set is
nonempty because $u=\|f\|_\infty h$ is feasible. Finite-dimensional
linear programming duality gives \eqref{eq:stencil-dual} and attainment.
The visit-measure formulation is the usual total-cost stopping
formulation; see \cite{Altman}. The argument above states its complete
finite-state realization in the present experiment.

For completeness, finite-horizon stopping gives the increasing
iteration $u_0=0$, $u_{n+1}=\max(0,Qu_n-f)$. Its values are bounded
by $\|f\|_\infty h$ and converge to the stopping value. The Bellman
map satisfies the componentwise estimate
$|B(a)-B(b)|\le Q|a-b|$. Hence any two fixed points differ in
absolute value by at most $Q^n$ times their original difference.
Transience proves uniqueness. Its value at $o$ equals
\eqref{eq:stencil-functional} by the visit representation.
\end{proof}

\begin{theorem}[Exact fixed-stencil occupation recovery]
\label{thm:finite-stencil}
For every target $x\in\R^2$, put $f_x(z)=g(x+tz)$, $z\in S$.
Then
\begin{equation}\label{eq:finite-stencil-inverse}
                         v(x)=\mathcal V(f_x).
\end{equation}
Thus precisely the same finite-dimensional functional recovers the
occupation at every point from a prior-bounded stencil. The positive
set, footprint, density, and containing component are not inputs.
No periodicity or boundary regularity is required.

The functional is defined also for data that are not the response of
any physical experiment. For arbitrary real vectors $f,\widetilde f$,
\begin{equation}\label{eq:stencil-stability}
 |\mathcal V(f)-\mathcal V(\widetilde f)|
 \le\sum_{z\in S}G_{oz}|f_z-\widetilde f_z|
 \le h_o\|f-\widetilde f\|_\infty.
\end{equation}
In particular, arbitrary measurement errors incur no factor equal to
an iteration horizon.
\end{theorem}
\begin{proof}
By the preceding lemma, $\mathcal V(f_x)$ is the supremum of
$\E_x[-\sum_{k<\tau}g(X_k)]$ over rules stopping by exit from
$x+tS$. Every such time has finite expectation. The martingale
identity used in Theorem~\ref{thm:global-occupation} gives
\[
 -\E_x\sum_{k<\tau}g(X_k)
                  =v(x)-\E_xv(X_\tau)\le v(x).
\]
If $v(x)=0$, stopping immediately attains equality. Otherwise $x$
belongs to one expanded component $P_C$. Every point of $P_C$ lies
within distance $D+\Delta\le Kt$ of $x$. Consequently the first exit
from its interior occurs no later than exit from $x+tS$ along the
compass walk. At that first exit occupation is zero, since one step
cannot reach another positive component. This admissible stopping rule
attains $v(x)$, proving \eqref{eq:finite-stencil-inverse}.

Other components may meet the edge of the stencil, and occupation
there may be positive. The inequality for every stopping rule and
the equality for the containing-component exit are precisely why no
physical zero boundary condition is needed.

For every feasible $m$, \eqref{eq:stencil-budget} bounds the change
of its objective by $\sum_zG_{oz}|f_z-\widetilde f_z|$. Taking maxima
in both directions proves \eqref{eq:stencil-stability}.
\end{proof}

\begin{corollary}[Local scalar stability]
\label{cor:stencil-local-stability}
Two admissible experiments with the same stated geometric bounds,
but possibly different configurations, footprints, and densities,
satisfy
\begin{equation}\label{eq:stencil-local-stability}
 |v(x)-\widetilde v(x)|
       \le h_o\max_{z\in S}|g(x+tz)-\widetilde g(x+tz)|.
\end{equation}
For a fixed experiment and every vector $a$,
\begin{equation}\label{eq:stencil-period-defect}
 |v(x+a)-v(x)|
       \le h_o\max_{z\in S}|g(x+a+tz)-g(x+tz)|.
\end{equation}
These are local estimates on fixed nominal apertures, not whole-field
error assumptions.
\end{corollary}
\begin{proof}
Apply the same functional at the two targets or experiments and use
\eqref{eq:stencil-stability}. The stencil matrix does not change under
translation.
\end{proof}

\subsection{Finite observations and numerical certificates}
The inverse gives both a finite sample law and a directly checkable
interval for every numerical output. If
$|f_z-\widehat f_z|\le\xi_z$ with $\xi_z\ge0$, then
\begin{equation}\label{eq:stencil-certified-interval}
 \mathcal V(\widehat f+\xi)\le\mathcal V(f)
                      \le\mathcal V(\widehat f-\xi).
\end{equation}
This follows because $m\ge0$ in every objective. Any feasible primal
vector supplies a lower bound and any feasible dual vector supplies
an upper bound in \eqref{eq:stencil-dual}. Thus finite rational
feasibility checks and a primal--dual gap certify the computed interval;
an unverified numerical optimizer status is unnecessary.

\begin{theorem}[A fixed-aperture finite occupation experiment]
\label{thm:stencil-sampling}
Fix $0<\varepsilon<1$ and $0<\delta<1$. At any nominal target $x$,
there is a procedure using at most $M$ distinct nominal sites and
\begin{equation}\label{eq:stencil-attempt-count}
 N\le 2M\left\lceil
       8H_S^2\varepsilon^{-2}\log\frac{4M}{\delta}
                         \right\rceil
\end{equation}
attempted collision bits, whose output has error at most
$\varepsilon$ with probability at least $1-\delta$. All sites lie
in $x+[-Kt,Kt]^2$, independently of the target accuracy. The result
is uniform over the unknown stationary density and convex footprint
in the exact class, without a quantitative boundary-mass assumption.
It remains valid conditionally at a target chosen from previous data.

There are $2M$ command-labelled batches, one forward and one reverse
per site; every batch repetition is included in $N$. Numerical
arithmetic is separate from these observation counts. If $x$ and $t$
are rational, all sites and both linear programs have rational data;
no progressively finer set of nominal positions is needed for this
occupation query.
\end{theorem}
\begin{proof}
At each site use $n$ fresh attempts for each mean, where $n$ is the
ceiling in \eqref{eq:stencil-attempt-count}, and let $\widehat f$
be the empirical forward-minus-reverse vector. Hoeffding's inequality
and a union bound give
\[
 \Prob\left\{\max_z|\widehat f_z-f_z|>
                      \frac{\varepsilon}{2H_S}\right\}
 \le4M\exp\left(-\frac{n\varepsilon^2}{8H_S^2}\right)
 \le\delta.
\]
Indeed it suffices to estimate each of the $2M$ means to
$\varepsilon/(4H_S)$. On this event
\eqref{eq:stencil-stability} bounds the error of
$\mathcal V(\widehat f)$ by $\varepsilon/2$. Compute its value to
certified error at most $\varepsilon/2$ using the finite rational
primal and dual programs, then project the result onto $[0,1]$.
Projection cannot increase the error because $v(x)\in[0,1]$.

Every recorded bit, including a solid start or a free miss, is an
attempt in the displayed count. Fresh preparations give exactly the
same bounds conditionally on previous observations. The stochastic
walk and stopping rule are computational objects and require no
additional physical flights. The rational coefficient and batch-count
bit lengths grow as $O(\log(1/\varepsilon)+\log\log(4M/\delta))$
after the rational target, step and fixed priors are specified.
\end{proof}

\begin{remark}[Aperture, interpretation, and comparison]
\label{rem:stencil-scope}
A bound on the absolute position of the footprint is needed only to
translate the nominal aperture into a fixed physical launch aperture:
forward and reverse starts lie in
$x+[-Kt,Kt]^2+A+t\overline B$. The number $M$ may be large; its
explicit dependence on the ratio $(D+\Delta)/t$ is not suppressed in
\eqref{eq:stencil-constants}. For a list of adaptive targets, allocate
conditional failure probabilities and add their batch costs. This
pointwise occupation theorem does not turn finitely many values into
a whole-field periodicity certificate.

The iteration of Lemma~\ref{lem:stopped} and
Corollary~\ref{cor:global-field-locality} remains valid. The present
finite program removes the horizon-dependent amplification for noisy,
not necessarily physically admissible, forcing. It does not replace
the one-sided local rare event of Proposition~\ref{prop:rare-query}:
estimating a signed occupation value and detecting a physical rare
collision are different statistical operations. In particular the
sharp stationary boundary power of
Corollary~\ref{cor:stationary-minimax} is unchanged.
\end{remark}
